\section{Grafos}



\subsection{MST --- Prim}

Adequado para grafos densos representados por matriz de adjacência.

\textbf{Complexidade:} $O(V^2)$.

\begin{lstlisting}
ll prim (const vector<vector<ll>>& graph) {
    int n = graph.size();
    vector<ll> minEdge(n, INF);
    vector<bool> used(n, false);
    ll totW = 0;
    minEdge[0] = 0;

    for (int i = 0; i < n; i++) {
        int curr = -1;
        for (int j = 0; j < n; j++)
            if (!used[j] && (curr == -1 || minEdge[j] < minEdge[curr])) curr = j;
        if (curr == -1 || minEdge[curr] == INF) return -1;
        used[curr] = true;
        totW += minEdge[curr];
        for (int k = 0; k < n; k++)
            if (!used[k] && graph[curr][k] < minEdge[k]) minEdge[k] = graph[curr][k];
    }
    return totW;
}
\end{lstlisting}



\subsection{Bellman-Ford}

\textbf{Complexidade:} $O(V E)$.

\begin{lstlisting}
using Edge = tuple<ll, int, int>;
pair<vector<ll>, bool> bellmanFord (int n, const vector<Edge>& edges, int origin) {
    vector<ll> dist(n, INF);
    dist[origin] = 0;
    for (int i = 0; i < n - 1; i++) {
        bool changed = false;
        for (auto [u, v, w] : edges) {
            if (dist[u] == INF) continue;
            if (dist[u] + w < dist[v]) {
                dist[v] = dist[u] + w;
                changed = true;
            }
        }
        if (!changed) break;
    }
    for (const Edge& edge : edges) {
        if (dist[u] == INF) continue;
        if (dist[u] + w < dist[v]) return {dist, true};
    }
    return {dist, false};
}
\end{lstlisting}


\subsection{Lowest Common Ancestor (Binary Lifting)}

\textbf{Pré-Processamento:} $O(N \log N)$.

\textbf{Consulta:} $O(\log N)$.

\begin{lstlisting}
const int MAXN = 2e5 + 5;
const int LOGN = 20;
vector<int> adj[MAXN];
int up[MAXN][LOGN];
int depth[MAXN];

void dfs_lca(int u, int p) {
    up[u][0] = p;
    for (int i = 1; i < LOGN; i++)
        up[u][i] = up[up[u][i - 1]][i - 1];
    for (int v : adj[u]) {
        if (v == p) continue;
            depth[v] = depth[u] + 1;
            dfs_lca(v, u);
    }
}
int lca(int u, int v) {
    if (depth[u] < depth[v]) swap(u, v);
    for (int i = LOGN - 1; i >= 0; i--) {
        if (depth[u] - (1 << i) >= depth[v]) 
            u = up[u][i];
    }
    if (u == v) return u;
    for (int i = LOGN - 1; i >= 0; i--) {
        if (up[u][i] != up[v][i]) {
            u = up[u][i];
            v = up[v][i];
        }
    }
    return up[u][0];
}
\end{lstlisting}


\subsection{Heavy-Light Decomposition (HLD)}

\textbf{Complexidade:} $O(\log^2 N)$.

\begin{lstlisting}
vector<int> adj[MAXN];
int sz[MAXN], depth[MAXN], parent[MAXN], heavy[MAXN], head[MAXN], pos[MAXN];
int cur_pos;

void dfs_sz(int u, int p, int d) {
    sz[u] = 1; depth[u] = d; parent[u] = p;
    heavy[u] = -1;
    int max_sz = 0;
    for (int v : adj[u]) {
        if (v != p) {
            dfs_sz(v, u, d + 1);
            sz[u] += sz[v];
            if (sz[v] > max_sz) { max_sz = sz[v]; heavy[u] = v; }
        }
    }
}
void dfs_hld(int u, int h) {
    head[u] = h;
    pos[u] = ++cur_pos;
    if (heavy[u] != -1) dfs_hld(heavy[u], h);
    for (int v : adj[u]) 
        if (v != parent[u] && v != heavy[u]) dfs_hld(v, v);
}
\end{lstlisting}


\subsection{Centroid Decomposition}

\textbf{Complexidade:} $O(N \log N)$.

\begin{lstlisting}
vector<set<int>> adj_c; // Grafo mutavel (arestas removidas)
vector<int> sz_c, parent_centroid;

void get_sz(int u, int p) {
    sz_c[u] = 1;
    for (int v : adj_c[u])
        if (v != p) { get_sz(v, u); sz_c[u] += sz_c[v]; }
}
int get_centroid(int u, int p, int total) {
    for (int v : adj_c[u])
        if (v != p && sz_c[v] > total / 2) return get_centroid(v, u, total);
    return u;
}
void decompose(int u, int p) {
    get_sz(u, p);
    int centroid = get_centroid(u, p, sz_c[u]);
    parent_centroid[centroid] = p;
    for (int v : adj_c[centroid]) {
        adj_c[v].erase(centroid);
        decompose(v, centroid);
    }
    adj_c[centroid].clear();
}
\end{lstlisting}


\subsection{Dinic's Algorithm (Max Flow)}

\textbf{Complexidade:} $O(V^2 E)$ (muito mais rápido na prática).

\begin{lstlisting}
struct Edge { int to; ll cap, flow; int rev; };
struct Dinic {
    int n;
    vector<vector<Edge>> adj;
    vector<int> level, ptr;
    Dinic(int n) : n(n), adj(n), level(n), ptr(n) {}
    void add_edge(int from, int to, ll cap) {
        adj[from].push_back({to, cap, 0, (int)adj[to].size()});
        adj[to].push_back({from, 0, 0, (int)adj[from].size() - 1});
    }
    bool bfs(int s, int t) {
        fill(level.begin(), level.end(), -1);
        level[s] = 0; queue<int> q; q.push(s);
        while (!q.empty()) {
            int v = q.front(); q.pop();
            for (auto& edge : adj[v]) {
                if (edge.cap - edge.flow > 0 && level[edge.to] == -1) {
                    level[edge.to] = level[v] + 1; q.push(edge.to);
                }
            }
        }
        return level[t] != -1;
    }
    ll dfs(int v, int t, ll pushed) {
        if (pushed == 0 || v == t) return pushed;
        for (int& cid = ptr[v]; cid < adj[v].size(); cid++) {
            auto& edge = adj[v][cid];
            int tr = edge.to;
            if (level[v] + 1 != level[tr] || edge.cap - edge.flow == 0) continue;
            ll tr_pushed = dfs(tr, t, min(pushed, edge.cap - edge.flow));
            if (tr_pushed == 0) continue;
            edge.flow += tr_pushed;
            adj[tr][edge.rev].flow -= tr_pushed;
            return tr_pushed;
        }
        return 0;
    }
    ll max_flow(int s, int t) {
        ll flow = 0;
        while (bfs(s, t)) {
            fill(all(ptr), 0);
            while (ll pushed = dfs(s, t, 1e18)) flow += pushed;
        }
        return flow;
    }
};
\end{lstlisting}


\subsection{Pontes em Grafos (Tarjan)}

\textbf{Complexidade:} $O(N + M)$.

\begin{lstlisting}
vector<int> adj_t[MAXN];
int tin[MAXN], low[MAXN], timer_t;
vector<pair<int, int>> bridges;

void dfs_bridge(int v, int p = -1) {
    tin[v] = low[v] = ++timer_t;
    for (int to : adj_t[v]) {
        if (to == p) continue;
        if (tin[to]) low[v] = min(low[v], tin[to]);
        else {
            dfs_bridge(to, v);
            low[v] = min(low[v], low[to]);
            if (low[to] > tin[v]) bridges.push_back({min(v, to), max(v, to)});
        }
    }
}
\end{lstlisting}


\subsection{Hopcroft-Karp (Maximum Bipartite Matching)}

\textbf{Complexidade:} $O(E \sqrt{V})$.

\begin{lstlisting}
struct HopcroftKarp {
    int n;
    vector<vector<int>> adj;
    vector<int> pair_u, pair_v, dist;
    HopcroftKarp(int n, int m) : n(n), adj(n + 1), pair_u(n + 1, 0), pair_v(m + 1, 0), dist(n + 1) {}
    void add_edge(int u, int v) { adj[u].push_back(v); }
    bool bfs() {
        queue<int> q;
        for (int u = 1; u <= n; u++) {
            if (pair_u[u] == 0) { dist[u] = 0; q.push(u); }
            else dist[u] = 1e9;
        }
        dist[0] = 1e9;
        while (!q.empty()) {
            int u = q.front(); q.pop();
            if (u != 0) {
                for (int v : adj[u]) {
                    if (dist[pair_v[v]] == 1e9) {
                        dist[pair_v[v]] = dist[u] + 1; q.push(pair_v[v]);
                    }
                }
            }
        }
        return dist[0] != 1e9;
    }
    bool dfs(int u) {
        if (u != 0) {
            for (int v : adj[u]) {
                if (dist[pair_v[v]] == dist[u] + 1 && dfs(pair_v[v])) {
                    pair_v[v] = u; pair_u[u] = v; return true;
                }
            }
            dist[u] = 1e9; return false;
        }
        return true;
    }
    int max_matching() {
        int matching = 0;
        while (bfs()) {
            for (int u = 1; u <= n; u++) if (pair_u[u] == 0 && dfs(u)) matching++;
        }
        return matching;
    }
};
\end{lstlisting}


\subsection{Min-Cost Max-Flow (MCMF)}

\textbf{Complexidade:} $O(E \sqrt{V})$.

\begin{lstlisting}
struct Edge {
    int to, cap, cost, flow, rev;
};

struct MCMF {
    int n;
    vector<vector<Edge>> adj;
    vector<ll> dist, potential;
    vector<int> parent_edge, parent_node;

    MCMF(int n) : n(n), adj(n), dist(n), potential(n, 0), parent_edge(n), parent_node(n) {}

    void add_edge(int from, int to, int cap, int cost) {
        adj[from].push_back({to, cap, cost, 0, (int)adj[to].size()});
        adj[to].push_back({from, 0, -cost, 0, (int)adj[to].size() - 1});
    }
    bool dijkstra(int s, int t) {
        fill(all(dist), 1e18);
        priority_queue<pair<ll, int>, vector<pair<ll, int>>, greater<>> pq;
        dist[s] = 0;
        pq.push({0, s});

        while (!pq.empty()) {
            auto [d, u] = pq.top(); pq.pop();
            if (d > dist[u]) continue;
            for (int i = 0; i < adj[u].size(); i++) {
                auto& edge = adj[u][i];
                if (edge.cap - edge.flow > 0 && dist[edge.to] > dist[u] + edge.cost + potential[u] - potential[edge.to]) {
                    dist[edge.to] = dist[u] + edge.cost + potential[u] - potential[edge.to];
                    parent_node[edge.to] = u;
                    parent_edge[edge.to] = i;
                    pq.push({dist[edge.to], edge.to});
                }
            }
        }
        for (int i = 0; i < n; i++) {
            if (dist[i] < 1e18) potential[i] += dist[i];
        }
        return dist[t] < 1e18;
    }
    pair<ll, ll> min_cost_max_flow(int s, int t) {
        ll flow = 0, cost = 0;
        while (dijkstra(s, t)) {
            ll push = 1e18;
            for (int v = t; v != s; v = parent_node[v]) {
                auto& edge = adj[parent_node[v]][parent_edge[v]];
                push = min(push, (ll)(edge.cap - edge.flow));
            }
            for (int v = t; v != s; v = parent_node[v]) {
                auto& edge = adj[parent_node[v]][parent_edge[v]];
                edge.flow += push;
                adj[v][edge.rev].flow -= push;
                cost += push * edge.cost;
            }
            flow += push;
        }
        return {flow, cost};
    }
};
\end{lstlisting}


\subsection{2-SAT Solver}

\textbf{Complexidade:} $O(N + M)$.

\begin{lstlisting}
struct TwoSAT {
    int n;
    vector<vector<int>> adj;
    vector<int> comp, tin, low, st;
    vector<bool> answer, on_stack;
    int timer, scc_cnt;

    TwoSAT(int n) : n(n), adj(2 * n), comp(2 * n, -1), tin(2 * n, 0), low(2 * n, 0), answer(n), on_stack(2 * n, false), timer(0), scc_cnt(0) {}

    void add_clause(int u, bool neg_u, int v, bool neg_v) {
        adj[2 * u + !neg_u].push_back(2 * v + neg_v);
        adj[2 * v + !neg_v].push_back(2 * u + neg_u);
    }
    void dfs(int u) {
        tin[u] = low[u] = ++timer;
        st.push_back(u); on_stack[u] = true;
        for (int v : adj[u]) {
            if (!tin[v]) {
                dfs(v);
                low[u] = min(low[u], low[v]);
            } 
            else if (on_stack[v])
                low[u] = min(low[u], tin[v]);
        }
        if (low[u] == tin[u]) {
            while (true) {
                int v = st.back(); st.pop_back();
                on_stack[v] = false;
                comp[v] = scc_cnt;
                if (u == v) break;
            }
            scc_cnt++;
        }
    }
    bool solve() {
        for (int i = 0; i < 2 * n; i++) if (!tin[i]) dfs(i);
        for (int i = 0; i < n; i++) {
            if (comp[2 * i] == comp[2 * i + 1]) return false;
            answer[i] = (comp[2 * i] > comp[2 * i + 1]);
        }
        return true;
    }
};
\end{lstlisting}